\documentclass{article}

\usepackage{booktabs}

\title{ELE 548: Branch Predictors}
\author{Calvin Higgins}

\begin{document}

\maketitle

\section{Tournament Branch Predictor}

\subsection{Design}

\subsubsection{Prediction}

The tournament branch predictor consists of three components:
\begin{enumerate}
    \item GAg predictor. The global branch history is mapped to a saturating counter via a table-lookup, and used to make a prediction.
    \item PAg predictor. The program counter is mapped to a local branch history via a table-lookup. Then, the local branch history is mapped to a saturating 
    counter via another table-lookup, and used to make a prediction. The second table is shared between all local branch histories.
    \item Bimodal choice predictor. The program counter is mapped to a saturating counter via a table-lookup, and used to predict whether the GAg or PAg 
    prediction should be used.
\end{enumerate}

\noindent
The size of the four tables, the number of bits in the global and local branch histories, and the number of bits in the saturating counters are left as 
compile-time parameters. The saturating counters are always initialized to weakly-taken. This implies that GAg is the default predictor during warmup.

\subsubsection{Update}

Each of the three components are updated independently:
\begin{enumerate}
    \item GAg predictor. The global branch history is extended with the actual branch result. The saturating counter is increased if the branch was taken, and 
    decreased otherwise.
    \item PAg predictor. The local branch history is extended with the actual branch result. The saturating counter is increased if the branch was taken, and
    decreased otherwise.
    \item Bimodal choice predictor. The update only occurs if the GAg and PAg predictions disagreed. In this case, the saturating counter is increased if GAg 
    was correct and decreased if PAg was correct.
\end{enumerate}

\subsection{Hyperparameter Sweep}
\label{sec:tbp_hp}

All combinations of compile-time parameters detailed in Table~\ref{tab:tbp_hp} were considered. Hyperparameter combinations that exceeded the strict hardware 
budget of 64KiB ($2^{19}$ bits) were rejected. Likewise, hyperparameter combinations that over-allocated the global or local saturating counter tables (i.e. 
when the history length is insufficient to index every entry), or did not come within 10\% of the hardware budget were rejected. The remaining 118 
hyperparameter combinations were simulated. 

\begin{table}[h]
    \centering
    \begin{tabular}{ll}
        \toprule
        \textbf{Hyperparameter} & \textbf{Values} \\
        \midrule
        Global history register bits       & $\{17, 18\}$ \\
        Global saturating counter table size & $\{2^{17}, 2^{18}\}$ \\
        Global saturating counter bits     & $\{2, 3\}$ \\
        Local history table size           & $\{2^{10}, 2^{12}, 2^{14}\}$ \\
        Local history bits                 & $\{10, 12, 14\}$ \\
        Local saturating counter table size & $\{2^{10}, 2^{12}, 2^{14}\}$ \\
        Local saturating counter bits      & $\{2, 3\}$ \\
        Choice saturating counter table size & $\{2^{10}, 2^{12}, 2^{14}\}$ \\
        Choice saturating counter bits     & $\{2, 3\}$ \\
        \bottomrule
    \end{tabular}
    \caption{Hyperparameter sweep for the tournament branch predictor.}
    \label{tab:tbp_hp}
\end{table}

\noindent
Due to the large search space, I only report MPKIs for the top three hyperparameter combinations. The results are provided in Table~\ref{tab:tbp_res}. TODO.

\subsubsection{Gshare}

The hyperparameter sweep detailed in Section~\ref{sec:tbp_hp} was repeated using Gshare instead of GAg. Again, due to the large search space, I only report 
MPKI for the top three hyperparameter combinations. The results are provided in Table~\ref{tab:tbp_gshare_res}. TODO.

\section{Hashed Perceptron Branch Predictor}

\subsection{Design}

Describe your final design and its configuration

\subsection{Hyperparameter Sweep}

 Report the MPKI results for each of the traces. Also report the
arithmetic and geometric mean MPKI of all traces. See Table 1.

\section{Custom Branch Predictor}

\subsection{Design}

Explain your final design.

\subsection{Hyperparameter Sweep}

Report the MPKI results for each of the traces. Also report the arithmetic and geometric
mean MPKI of all traces. Compare with TAGE-SC-L and your Problem 2 results


\newpage
\appendix

\begin{table}[h]
    \centering
    \begin{tabular}{llll}
        \toprule
        \textbf{Trace} & \textbf{TODO MPKI} & \textbf{TODO MPKI} & \textbf{TODO MPKI} \\
        \midrule
        SHORT-FP-1  &  &  & \\
        SHORT-FP-2  &  &  & \\
        SHORT-FP-3  &  &  & \\
        SHORT-FP-4  &  &  & \\
        SHORT-FP-5  &  &  & \\
        SHORT-INT-1  &  &  & \\
        SHORT-INT-2  &  &  & \\
        SHORT-INT-3  &  &  & \\
        SHORT-INT-4  &  &  & \\
        SHORT-INT-5  &  &  & \\
        SHORT-MM-1  &  &  & \\
        SHORT-MM-2  &  &  & \\
        SHORT-MM-3  &  &  & \\
        SHORT-MM-4  &  &  & \\
        SHORT-MM-5  &  &  & \\
        SHORT-SERV-1  &  &  & \\
        SHORT-SERV-2  &  &  & \\
        SHORT-SERV-3  &  &  & \\
        SHORT-SERV-4  &  &  & \\
        SHORT-SERV-5  &  &  & \\
        LONG-SPEC2K6-00  &  &  & \\
        LONG-SPEC2K6-01  &  &  & \\
        LONG-SPEC2K6-02  &  &  & \\
        LONG-SPEC2K6-03  &  &  & \\
        LONG-SPEC2K6-04  &  &  & \\
        LONG-SPEC2K6-05  &  &  & \\
        LONG-SPEC2K6-06  &  &  & \\
        LONG-SPEC2K6-07  &  &  & \\
        LONG-SPEC2K6-08  &  &  & \\
        LONG-SPEC2K6-09  &  &  & \\
        LONG-SPEC2K6-10  &  &  & \\
        LONG-SPEC2K6-11  &  &  & \\
        LONG-SPEC2K6-12  &  &  & \\
        LONG-SPEC2K6-13  &  &  & \\
        LONG-SPEC2K6-14  &  &  & \\
        LONG-SPEC2K6-15  &  &  & \\
        LONG-SPEC2K6-16  &  &  & \\
        LONG-SPEC2K6-17  &  &  & \\
        LONG-SPEC2K6-18  &  &  & \\
        LONG-SPEC2K6-19  &  &  & \\
        \midrule
        \textbf{Arithmetic Mean}  &  &  & \\
        \textbf{Geometric Mean}  &  &  & \\
        \bottomrule
    \end{tabular}
    \caption{Results for top-three tournament branch predictor combinations.}
    \label{tab:tbp_res}
\end{table}

\begin{table}[h]
    \centering
    \begin{tabular}{llll}
        \toprule
        \textbf{Trace} & \textbf{TODO MPKI} & \textbf{TODO MPKI} & \textbf{TODO MPKI} \\
        \midrule
        SHORT-FP-1  &  &  & \\
        SHORT-FP-2  &  &  & \\
        SHORT-FP-3  &  &  & \\
        SHORT-FP-4  &  &  & \\
        SHORT-FP-5  &  &  & \\
        SHORT-INT-1  &  &  & \\
        SHORT-INT-2  &  &  & \\
        SHORT-INT-3  &  &  & \\
        SHORT-INT-4  &  &  & \\
        SHORT-INT-5  &  &  & \\
        SHORT-MM-1  &  &  & \\
        SHORT-MM-2  &  &  & \\
        SHORT-MM-3  &  &  & \\
        SHORT-MM-4  &  &  & \\
        SHORT-MM-5  &  &  & \\
        SHORT-SERV-1  &  &  & \\
        SHORT-SERV-2  &  &  & \\
        SHORT-SERV-3  &  &  & \\
        SHORT-SERV-4  &  &  & \\
        SHORT-SERV-5  &  &  & \\
        LONG-SPEC2K6-00  &  &  & \\
        LONG-SPEC2K6-01  &  &  & \\
        LONG-SPEC2K6-02  &  &  & \\
        LONG-SPEC2K6-03  &  &  & \\
        LONG-SPEC2K6-04  &  &  & \\
        LONG-SPEC2K6-05  &  &  & \\
        LONG-SPEC2K6-06  &  &  & \\
        LONG-SPEC2K6-07  &  &  & \\
        LONG-SPEC2K6-08  &  &  & \\
        LONG-SPEC2K6-09  &  &  & \\
        LONG-SPEC2K6-10  &  &  & \\
        LONG-SPEC2K6-11  &  &  & \\
        LONG-SPEC2K6-12  &  &  & \\
        LONG-SPEC2K6-13  &  &  & \\
        LONG-SPEC2K6-14  &  &  & \\
        LONG-SPEC2K6-15  &  &  & \\
        LONG-SPEC2K6-16  &  &  & \\
        LONG-SPEC2K6-17  &  &  & \\
        LONG-SPEC2K6-18  &  &  & \\
        LONG-SPEC2K6-19  &  &  & \\
        \midrule
        \textbf{Arithmetic Mean}  &  &  & \\
        \textbf{Geometric Mean}  &  &  & \\
        \bottomrule
    \end{tabular}
    \caption{Results for top-three tournament branch predictor combinations with Gshare.}
    \label{tab:tbp_gshare_res}
\end{table}

\end{document}